\scanpage{50}
\mspage{\handwritten{6}}
% The typescript is upside down in the scan and crossed by long diagonal pen strokes.
\deleted{Pour la première assertion, on se ramène encore, par descente et en utilisant \deleted{1}\handwritten{4}.\handwritten{3}, au cas où $k$ est séparablement clos. Procédant comme dans la démonstration de \deleted{1}\handwritten{4}.6 , on déduit de \deleted{1}\handwritten{4}.5 et de l'hypothèse sur $L$ que $\underline{L}$ \deleted{provient} est isomorphe à l'image inverse d'un Module birigidifié $\underline{M}$ sur $A\mathbin{\mathrm{x}}Q$, et pour la première assertion de \deleted{1}\handwritten{4}.7, on est donc ramené au cas $A=P$. Or $\underline L$ définit un morphisme de schémas ponctués}
\[
 \deleted{f\colon Q\handwritten{\longrightarrow}\underline{\mathrm{P}\deleted{\unsure{i|j}}\mathrm{ic}}_{P/k}\ ,}
\]
\deleted{qui se factorise nécessairement par le schéma abélien dual $P'$ de $P$. Soit $\underline{L}'$ l'image inverse du faisceau de Weil sur $P\mathbin{\mathrm{x}}P'$ par \deleted{$\mathrm{id}_{P}\mathbin{\mathrm{x}}\mathrm{id}$}}
\[
 \deleted{\deleted{\unsure{p|P}}\mathrm{id}_{P}\mathbin{\mathrm{x}}f\colon P\mathbin{\mathrm{x}}Q\handwritten{\longrightarrow}P\mathbin{\mathrm{x}}P'\ ,}
\]
\deleted{c'est un Module birigidifié qui, par définition de $f$, ne diffère de $\underline L$ que modulo l'image inverse d'un Module inversible $\underline M$ sur $P$. Comme $\underline L'$ et $\underline L$ ont une restriction triviale à $\deleted{\unsure{E|B}}P\mathbin{\mathrm{x}}e_{Q}$ , on en conclut que $\underline M$ est \deleted{\unsure{trivaal|trivaial}}trivial, donc que $\underline L$ est isomorphe à $\underline L'$. Or nous avons signalé dans VII 2.9.4 \deleted{\unsure{p|q}}que sur un produit de deux schémas abéliens, les correspondances divisorielles provenaient de biextensions. Il en résulte que $\underline L'$ , \deleted{\unsure{et en|et ea}} ou \deleted{\unsure{p|pe}}ce qui revient au même, $\underline L$ , provient d'une biextension.}

\deleted{Revenant au cas général ($P$ extension de $A$ par $L$), \deleted{\unsure{et|él}}il reste à prouver que (\deleted{1}\handwritten{4}.7.1) est bien un isomorphisme. Mais cela résulte de la suite exacte des $\mathrm{B}{\unsure{\bar{\mathrm{i}}|\mathrm{i}}}\mathrm{ext}^{i}$ ($i=0,1$) associée à la suite exacte}
\[
 \deleted{0\handwritten{\longrightarrow}L\handwritten{\longrightarrow}P\handwritten{\longrightarrow}A\handwritten{\longrightarrow}0\ ,}
\]
\deleted{et des relations}
\[
 \deleted{\operatorname{Biext}^{\unsure{o|0}}(L,Q;\underline G_{m})=0\ ,\quad
 \operatorname{Biext}^{1}(L\overset{\deleted{\mathrm{x}}}{,}Q;\underline G_{m})=0\ ,}
\]
\deleted{dont la première est un cas \deleted{\unsure{partioulier|partiwulier}}particulier de l'assertion de rigidité faite dans \deleted{1}\handwritten{4}.\handwritten{3} , et la deuxième n'est autre que le premier cas de \deleted{1}\handwritten{4}.6.}

\deleted{\uline{Corollaire} \deleted{1}\handwritten{4}.8. Sous les conditions de \deleted{1}\handwritten{4}.7, supposons que $Q$ soit extension d'un groupe algébrique $B$ (\unsure{nécesairement|nécessairement} lisse et connexe) par un}
% A short vertical pen stroke brackets the two lines of the corollary.
% archive sheet number: 49
